\documentclass[9pt, aspectratio=169]{beamer}

% --- CONFIGURATION ROBUSTE (MacTeX / Overleaf) ---
\usepackage[utf8]{inputenc}
\usepackage[T1]{fontenc}
\usepackage{lmodern}
\usepackage[french]{babel}
% -------------------------------------------------

% --- THÈME ---
\usetheme{Warsaw}
\usecolortheme{default}
\setbeamertemplate{navigation symbols}{}
\setbeamertemplate{headline}{}
\setbeamertemplate{caption}[numbered]

% --- PACKAGES ---
\usepackage{amsmath}
\usepackage{amsfonts}
\usepackage{amssymb}
\usepackage{tikz}
\usetikzlibrary{shapes, arrows, positioning, calc, shadows, trees, fit, backgrounds, decorations.pathreplacing}
\usepackage{booktabs}
\usepackage{csquotes}
\usepackage{ragged2e}
\usepackage{caption}
\usepackage{hyperref}

% --- CONFIGURATION DES LIENS ---
\hypersetup{
    colorlinks=true,
    linkcolor=black,
    urlcolor=blue,
    citecolor=blue
}

% --- BIBLIOGRAPHIE (Mode Compatible BibTeX pour MacTeXnicle) ---
% On utilise 'bibtex' au lieu de 'biber' pour assurer la compatibilité
% avec les éditeurs configurés par défaut.
\usepackage[style=authoryear, backend=bibtex]{biblatex}

% --- CRÉATION DU FICHIER BIB ---
\begin{filecontents*}{ref_doctorat.bib}
@article{Weber2007The,
  author = {Weber, Max},
  title = {The Concept of Bureaucracy},
  journal = {International Journal of Social Science Studies},
  volume = {7},
  number = {2},
  year = {2019}
}
@book{Graeber2015,
  author = {Graeber, David},
  title = {Bureaucratie. L'utopie des règles},
  publisher = {Les Liens qui libèrent},
  year = {2015},
  address = {Paris}
}
@book{Zuboff2023The,
  author = {Zuboff, Shoshana},
  title = {The Age of Surveillance Capitalism},
  year = {2019},
  publisher = {PublicAffairs}
}
@misc{Biosphere2015,
  author = {{Biosphere Code}},
  title = {The Biosphere Code},
  year = {2015},
  note = {Manifeste pour une algorithmique responsable}
}
@article{Neuer2024Le,
  author = {Neuer, Laurence},
  title = {Le droit face aux algorithmes},
  journal = {Le Point},
  year = {2023}
}
@article{Doran2002What,
  author = {Doran, Derek and Schulz, Sarah and Besold, Tarek R.},
  title = {What Does Explainable AI Really Mean?},
  journal = {arXiv preprint arXiv:1710.00794},
  year = {2017}
}
@article{Burkart2021A,
  author = {Burkart, Nadia and Huber, Marco F.},
  title = {A Survey on the Explainability of Supervised Machine Learning},
  journal = {Journal of Artificial Intelligence Research},
  volume = {70},
  pages = {245--317},
  year = {2021}
}
@article{Ko2020A,
  author = {Ko, Hyeyoung and Lee, Suyeon and Park, Yoonseo and Choi, Anna},
  title = {A Survey of Recommendation Systems},
  journal = {Electronics},
  volume = {11},
  number = {1},
  year = {2022}
}
@article{Koren2022Advances,
  author = {Koren, Yehuda and Bell, Robert},
  title = {Advances in Collaborative Filtering},
  journal = {Recommender Systems Handbook},
  pages = {145--186},
  year = {2009},
  publisher = {Springer}
}
@book{Pariser2011The,
  author = {Pariser, Eli},
  title = {The Filter Bubble: What the Internet Is Hiding from You},
  year = {2011},
  publisher = {Penguin Press},
  address = {New York}
}
\end{filecontents*}

\addbibresource{ref_doctorat.bib}

% --- INFOS TITRE ---
\title{Le Léviathan Algorithmique}
\subtitle{Pouvoir, opacité et responsabilité à l'ère de l'IA}
\author{Dominique Loyer}
\institute{Université du Québec à Montréal (UQAM) \\ Doctorat en Informatique Cognitive}
\date{26 novembre 2025}

\begin{document}

% --- SLIDE 1 : Titre ---
\begin{frame}
    \titlepage
    \begin{center}
        \tiny Examen de synthèse doctoral (DIC-9401)
    \end{center}
\end{frame}

% --- SLIDE 2 : Plan ---
\begin{frame}{Plan de l'examen : 4 axes de réflexion}
    Ce rapport s'articule autour de quatre questionnements fondamentaux, soutenus par la littérature de référence.
    
    \vspace{0.3cm}
    \textbf{1. Bureaucratie Algorithmique}
    \begin{itemize}
        \item \textit{Problématique :} Lien entre personnalisation et phénomène bureaucratique.
        \item \textit{Sources :} \cite{Weber2007The}, \cite{Graeber2015}.
    \end{itemize}
    
    \textbf{2. Code de Déontologie}
    \begin{itemize}
        \item \textit{Problématique :} Responsabilité humaine et formulation d'un code éthique.
        \item \textit{Sources :} \cite{Biosphere2015}, \cite{Neuer2024Le}.
    \end{itemize}
    
    \textbf{3. Interprétabilité des Modèles}
    \begin{itemize}
        \item \textit{Problématique :} Définition, évaluation et enjeux des méthodes XAI.
        \item \textit{Sources :} \cite{Doran2002What}, \cite{Burkart2021A}.
    \end{itemize}
    
    \textbf{4. Systèmes de Recommandation}
    \begin{itemize}
        \item \textit{Problématique :} Fonctionnement technique et comparaison d'approches.
        \item \textit{Sources :} \cite{Ko2020A}, \cite{Koren2022Advances}, \cite{Pariser2011The}.
    \end{itemize}
\end{frame}

% ================= SECTION 1 : BUREAUCRATIE =================
\section{Bureaucratie Algorithmique}

\begin{frame}{Section 1 : Bureaucratie algorithmique}
    \begin{center}
        \Large \textbf{La rationalisation numérique du monde}
    \end{center}
    \vspace{1cm}
    \begin{quote}
        « Quel lien pourriez-vous faire entre les modalités de fonctionnement de l'algorithme de personnalisation et le phénomène bureaucratique ? »
    \end{quote}
\end{frame}

% --- FIGURE VECTORIELLE WEBER ---
\begin{frame}{Illustration : la hiérarchie wébérienne}
    \begin{figure}
        \centering
        \begin{tikzpicture}[
            node distance=1.5cm,
            every node/.style={rectangle, rounded corners, draw=black!80, fill=blue!10, align=center, font=\scriptsize, minimum width=2.0cm, drop shadow},
            arrow/.style={->, >=stealth, thick, color=black!70}
        ]
            \node (authority) {\textbf{Règles Abstraites}\\(Code = Loi)};
            \node (hierarchy) [below left=1cm and 0.2cm of authority] {Hiérarchie\\Pyramidale};
            \node (impersonal) [below right=1cm and 0.2cm of authority] {Impersonnalité\\(Standardisation)};
            \node (efficiency) [below=2.5cm of authority, fill=green!10] {\textbf{Rationalité}\\(Optimisation)};
            
            \draw[arrow] (authority) -- (hierarchy);
            \draw[arrow] (authority) -- (impersonal);
            \draw[arrow] (hierarchy) -- (efficiency);
            \draw[arrow] (impersonal) -- (efficiency);
            
            \node [draw=none, fill=none, right=0.5cm of efficiency, text width=3cm, font=\tiny\itshape] {Objectif : Efficacité Instrumentale};
        \end{tikzpicture}
        \caption{Modèle bureaucratique de Weber appliqué à l'IA (Généré par GEMINI PRO 3.0)}
    \end{figure}
\end{frame}

\begin{frame}{La vision wébérienne}
    Selon \cite{Weber2007The}, la bureaucratie se définit par :
    \begin{enumerate}
        \item \textbf{Règles abstraites et impersonnelles :} L'algorithme applique des règles de code sans considération pour l'individu en tant que sujet moral.
        \item \textbf{Rationalité instrumentale :} Optimisation de moyens en vue d'une fin (le clic, l'engagement) sans questionnement sur la valeur de la fin.
        \item \textbf{Opacité et Secret :} Le « secret de fonction » bureaucratique devient le « secret industriel » ou l'opacité technique de la boîte noire.
    \end{enumerate}
\end{frame}

\begin{frame}{L'algorithme comme bureaucrate parfait}
    L'algorithme de personnalisation réalise le rêve bureaucratique ultime :
    \vspace{0.5cm}
    \begin{itemize}
        \item \textbf{Application instantanée des règles :} Aucune latence entre la décision et l'exécution.
        \item \textbf{Absence d'états d'âme :} La suppression totale de l'empathie humaine, vue comme un biais dans la logique bureaucratique pure.
        \item \textbf{Traitement de masse standardisé :} Malgré l'illusion de personnalisation, chaque individu est traité selon des catégories statistiques pré-établies.
    \end{itemize}
\end{frame}

\begin{frame}{David Graeber et l'utopie des règles}
    L'analyse de \cite{Graeber2015} éclaire la nature contraignante de l'algorithme :
    \begin{itemize}
        \item \textbf{La violence structurelle :} L'impossibilité de négocier avec la machine. Face à une recommandation ou un refus algorithmique, il n'y a pas de recours. L'interface utilisateur agit comme un guichet fermé.
        \item \textbf{La zone morte de l'imagination :} La bureaucratie (et l'algorithme) remplace l'imagination humaine et l'empathie par des procédures rigides. L'algorithme ne peut pas imaginer des situations hors de ses données d'entraînement.
    \end{itemize}
\end{frame}

\begin{frame}{Le paradoxe de la personnalisation bureaucratique}
    Comment un système \textit{personnalisé} peut-il être \textit{bureaucratique} ?
    \begin{itemize}
        \item La personnalisation n'est pas une rencontre humaine, mais une \textbf{catégorisation statistique}.
        \item L'individu est réduit à un vecteur de données (le « dividu »).
        \item La gestion par algorithme est une administration des comportements à grande échelle, répondant parfaitement à la définition de la gestion bureaucratique moderne : gérer des populations par des normes.
    \end{itemize}
\end{frame}

\begin{frame}{Le capitalisme de surveillance}
    \cite{Zuboff2023The} fournit le cadre économique de cette bureaucratie :
    \begin{itemize}
        \item L'algorithme n'est pas seulement un outil de gestion, c'est un outil d'\textbf{extraction}.
        \item La rationalisation (Weber) est mise au service de la prédiction et de la modification comportementale à des fins lucratives.
        \item Le "surplus comportemental" est la matière première de cette nouvelle industrie, traitée par la bureaucratie algorithmique.
    \end{itemize}
\end{frame}

\begin{frame}{La "rationalité" du profilage}
    Le profilage algorithmique est une forme de rationalisation extrême :
    \begin{itemize}
        \item \textbf{Standardisation de l'humain :} Réduire la complexité d'une personne à un ensemble de « caractéristiques » (features) exploitables par le modèle.
        \item \textbf{Déqualification du jugement :} Le jugement humain est remplacé par un score de probabilité (score de crédit, score de recommandation).
        \item \textbf{Automatisation de la suspicion :} Le système est conçu pour trier, classer et évaluer en permanence, souvent de manière préemptive.
    \end{itemize}
\end{frame}

\begin{frame}{Synthèse Q1 : la machine comme guichet}
    L'algorithme de personnalisation est l'aboutissement de la bureaucratie :
    \begin{itemize}
        \item Il applique une \textbf{rationalité instrumentale} (Weber) sans empathie.
        \item Il impose une \textbf{violence structurelle} par son inflexibilité (Graeber).
        \item Il transforme l'expérience humaine en données extractibles pour un marché (Zuboff).
    \end{itemize}
    \vspace{0.3cm}
    La "personnalisation" est l'interface aimable d'un appareil bureaucratique de gestion et de contrôle impersonnel.
\end{frame}

\begin{frame}{Conséquences : la cage d'acier numérique}
    Weber parlait de « cage d'acier » pour la rationalité capitaliste. L'IA renforce cette cage :
    \begin{itemize}
        \item \textbf{Inflexibilité :} Le code est loi (Lessig). Une fois déployée, la règle s'applique sans discernement.
        \item \textbf{Irrésponsabilité :} "C'est la faute de l'ordinateur". La chaîne de responsabilité est diluée entre le développeur, les données et l'utilisateur.
        \item \textbf{Perte de sens :} L'efficacité (le "comment") prime systématiquement sur la finalité sociale (le "pourquoi").
    \end{itemize}
\end{frame}

\begin{frame}{Le wébérisme numérique : formalisation}
    \textbf{Concepts clés sur la bureaucratie algorithmique} :
    
    \vspace{0.3cm}
    \begin{itemize}
        \item \textbf{Administration isocratique :} L'IA permet une administration où l'explicabilité algorithmique augmente la compétence individuelle du personnel et des citoyens
        \item \textbf{Gestion de la complexité :} Les algorithmes traitent la complexité institutionnelle plus efficacement que les humains seuls
        \item \textbf{Bureaucratie hybride :} Mélange d'agence humaine et non-humaine dans la prise de décision
    \end{itemize}
    
    \vspace{0.3cm}
    \textit{Paradoxe :} La bureaucratie algorithmique peut être \textbf{impartiale} dans son application tout en étant \textbf{sensible au contexte} de façon prédictive.
\end{frame}

% ================= SECTION 2 : DÉONTOLOGIE =================
\section{Code de Déontologie}

\begin{frame}{Section 2 : Code de déontologie}
    \begin{center}
        \Large \textbf{Responsabilité et éthique du codeur}
    \end{center}
    \vspace{1cm}
    \begin{quote}
        « Comment vous y prendriez-vous pour formuler un code de déontologie des codeurs... ? »
    \end{quote}
\end{frame}

\begin{frame}{Inspiration : The Biosphere Code}
    Le \textit{Biosphere Code} \parencite{Biosphere2015} propose un manifeste liant technologie et écologie.
    \begin{itemize}
        \item Nous devons adapter cette approche à l'écologie \textit{attentionnelle} et \textit{sociale} des utilisateurs.
        \item L'algorithme n'est pas neutre ; il habite et modifie notre environnement mental comme une industrie modifie l'environnement naturel.
        \item Le codeur a une responsabilité civique qui dépasse son contrat de travail.
    \end{itemize}
\end{frame}

\begin{frame}{Le contexte juridique et social}
    \cite{Neuer2024Le} souligne l'urgence d'un cadre :
    \begin{itemize}
        \item Le vide juridique actuel laisse les développeurs seuls face à des dilemmes éthiques complexes.
        \item La pression organisationnelle pousse à la performance (KPIs, engagement) au détriment de l'éthique.
        \item Il est nécessaire de protéger le lanceur d'alerte interne qui signale des dérives algorithmiques.
    \end{itemize}
\end{frame}

% --- FIGURE VECTORIELLE PILIERS ---
\begin{frame}{Illustration : piliers du code}
    \begin{figure}
        \centering
        \begin{tikzpicture}[scale=0.8, every node/.style={transform shape}]
            % Sol
            \fill[gray!30] (-3,0) rectangle (3,0.3);
            \node[below] at (0,0) {\textit{Confiance des utilisateurs}};
            % Piliers
            \foreach \x/\c/\t in {-2/blue!30/Transparence, 0/green!30/Loyauté, 2/red!30/Responsabilité} {
                \draw[fill=\c] (\x-0.8, 0.3) rectangle (\x+0.8, 3.5);
                \node[rotate=90] at (\x, 1.9) {\textbf{\t}};
            }
            % Toit
            \draw[fill=black!70] (-3.2, 3.5) -- (3.2, 3.5) -- (0, 5) -- cycle;
            \node[white, font=\bfseries] at (0, 4) {DÉONTOLOGIE};
        \end{tikzpicture}
        \caption{Les trois piliers d'une déontologie algorithmique (Généré par GEMINI PRO 3.0)}
    \end{figure}
\end{frame}

\begin{frame}{Proposition de code : article 1 - transparence}
    \textbf{« Je ne coderai pas de boîtes noires destinées à manipuler. »}
    \vspace{0.5cm}
    \begin{itemize}
        \item Le développeur doit s'assurer que les critères d'optimisation sont explicables et documentés.
        \item Refus des « motifs obscurs » (dark patterns) dans la conception des systèmes de recommandation.
        \item L'utilisateur doit savoir pourquoi un contenu lui est recommandé.
    \end{itemize}
\end{frame}

\begin{frame}{Proposition de code : article 2 - loyauté}
    \textbf{« Mon algorithme servira l'intérêt de l'utilisateur, pas seulement celui de la plateforme. »}
    \vspace{0.5cm}
    \begin{itemize}
        \item Prioriser la pertinence réelle et le bien-être sur le temps passé (engagement addictif).
        \item Ne pas exploiter les vulnérabilités psychologiques des utilisateurs.
        \item \cite{Biosphere2015} rappelle la responsabilité envers l'écosystème global : un algorithme ne doit pas polluer l'espace public.
    \end{itemize}
\end{frame}

\begin{frame}{Proposition de code : article 3 - responsabilité}
    \textbf{« J'assume la responsabilité des biais de mon modèle. »}
    \vspace{0.5cm}
    \begin{itemize}
        \item Test rigoureux des biais discriminatoires (genre, race, classe) avant tout déploiement.
        \item Mise en place de mécanismes de monitoring continu post-déploiement.
        \item Création de voies de recours pour les utilisateurs s'estimant lésés par une décision automatique.
    \end{itemize}
\end{frame}

\begin{frame}{Processus d'élaboration (1/2) : consultation}
    Un code de déontologie ne peut être unilatéral.
    \begin{itemize}
        \item \textbf{Multidisciplinaire :} Doit impliquer des ingénieurs, mais aussi des juristes, sociologues, psychologues et représentants des utilisateurs.
        \item \textbf{Itératif :} Le code doit être un document vivant, capable de s'adapter aux nouvelles avancées technologiques (IA générative, etc.).
        \item \textbf{Public :} Le processus de consultation doit être transparent pour garantir l'adhésion.
    \end{itemize}
\end{frame}

\begin{frame}{Processus d'élaboration (2/2) : application}
    Pour être efficace, le code doit être contraignant.
    \begin{itemize}
        \item \textbf{Intégration au flux de travail (workflow) :} L'évaluation éthique (Ethical Impact Assessment) doit être une étape obligatoire du cycle de développement, comme les tests unitaires.
        \item \textbf{Cadre légal :} Le code doit s'appuyer sur des régulations émergentes (ex: AI Act européen) pour avoir une portée légale \parencite{Neuer2024Le}.
        \item \textbf{Protection :} Mécanismes de protection des ingénieurs qui refusent une tâche contraire à l'éthique (droit de retrait).
    \end{itemize}
\end{frame}

\begin{frame}{Cadres réglementaires émergents}
    \textbf{Principes internationaux pour une IA responsable} :
    
    \vspace{0.3cm}
    \textbf{1. Recommandation UNESCO sur l'éthique de l'IA (2021)} :
    \begin{itemize}
        \item Approche centrée sur l'humain, droits humains, diversité culturelle
        \item Transparence, responsabilité, contrôlabilité par l'humain
    \end{itemize}
    
    \textbf{2. Groupe d'experts de haut niveau de l'UE (2019)} :
    \begin{itemize}
        \item \textit{Quatre principes éthiques :} Respect de l'autonomie humaine, prévention des dommages, équité, explicabilité
        \item \textit{Sept exigences clés :} Agence humaine, robustesse technique, vie privée, transparence, diversité, bien-être sociétal, responsabilité
    \end{itemize}
    
    \textbf{3. Principes ACM/IEEE (2022-2023)} :
    \begin{itemize}
        \item Auditabilité, contestabilité, transparence, sécurité
        \item "Humain dans la boucle" obligatoire pour systèmes à haut risque
    \end{itemize}
\end{frame}

\begin{frame}{Redevabilité algorithmique : cadre opérationnel}
    \textbf{Trois dimensions de la redevabilité algorithmique} :
    
    \vspace{0.3cm}
    \begin{enumerate}
        \item \textbf{Responsabilité} : 
        \begin{itemize}
            \item Identification claire des acteurs responsables (développeurs, déployeurs, régulateurs)
            \item Chaîne de responsabilité sur tout le cycle de vie de l'IA
        \end{itemize}
        
        \item \textbf{Transparence} : 
        \begin{itemize}
            \item Systèmes explicables et auditables
            \item Détection et correction des erreurs
        \end{itemize}
        
        \item \textbf{Recours} : 
        \begin{itemize}
            \item Voies d'appel pour les individus lésés
            \item Mécanismes de contestation des décisions automatisées
        \end{itemize}
    \end{enumerate}
    
    \vspace{0.3cm}
    \textbf{Outils pratiques :} Évaluations d'impact algorithmique (ÉIA), audits indépendants, documentation du cycle de vie.
\end{frame}

\begin{frame}{Le "serment d'Hippocrate" du codeur}
    L'objectif final est de passer d'une culture de "Move Fast and Break Things" à une culture de la responsabilité.
    \begin{itemize}
        \item Le principe fondamental : \textbf{"Primum non nocere" (D'abord, ne pas nuire)}.
        \item La responsabilité du codeur n'est pas seulement technique (le code compile-t-il ?) mais aussi civique (le code est-il juste ?).
        \item Cela implique une formation en éthique obligatoire dans les cursus d'ingénierie.
    \end{itemize}
\end{frame}

\begin{frame}{Mise en \oe uvre organisationnelle}
    Un code n'est rien sans structure pour le soutenir :
    \begin{itemize}
        \item Création de comités d'éthique indépendants au sein des entreprises tech.
        \item Institution d'un droit de retrait pour un ingénieur refusant de coder une fonctionnalité contraire au code.
        \item Formation continue en éthique des données, pas seulement en technique.
    \end{itemize}
\end{frame}

% ================= SECTION 3 : INTERPRÉTABILITÉ =================
\section{Interprétabilité des Modèles}

\begin{frame}{Section 3 : Interprétabilité (XAI)}
    \begin{center}
        \Large \textbf{Ouvrir la boîte noire}
    \end{center}
    \vspace{1cm}
    \begin{quote}
        « Comment définit-on et évalue-t-on l'interprétabilité... ? »
    \end{quote}
\end{frame}

\begin{frame}{Définitions fondamentales}
    Selon \cite{Doran2002What}, il faut distinguer :
    \begin{itemize}
        \item \textbf{Interprétabilité (Interpretability) :} La capacité d'un modèle à être compris dans sa globalité par un humain. C'est une propriété intrinsèque (ex: régression linéaire, arbre de décision simple).
        \item \textbf{Explicabilité (Explainability) :} La capacité d'un système opaque à fournir des justifications a posteriori pour ses décisions (ex: cartes de chaleur pour un réseau de neurones).
        \item \textbf{Transparence :} Accès technique au code source et aux données d'entraînement.
    \end{itemize}
\end{frame}

\begin{frame}{Pourquoi est-ce crucial ?}
    \cite{Burkart2021A} identifient les enjeux majeurs :
    \begin{enumerate}
        \item \textbf{Confiance (Trust) :} L'utilisateur n'adopte pas durablement une technologie qu'il ne comprend pas ou qui lui semble arbitraire.
        \item \textbf{Débogage et Amélioration :} Pour l'ingénieur, comprendre pourquoi le modèle échoue est essentiel pour le corriger.
        \item \textbf{Conformité Légale :} RGPD (Droit à l'explication des décisions automatisées), AI Act européen.
        \item \textbf{Éthique :} Détection des biais cachés qui seraient invisibles par la seule analyse de la performance globale.
    \end{enumerate}
\end{frame}

\begin{frame}{Taxonomie des méthodes (1/2)}
    \textbf{1. Modèles intrinsèquement interprétables}
    \begin{itemize}
        \item \textbf{Régression Linéaire / Logistique :} Les poids associés à chaque variable sont lisibles directement.
        \item \textbf{Arbres de décision :} Suivre le chemin de la racine à la feuille est un processus logique humain.
        \item \textit{Avantage :} Transparence totale et confiance élevée.
        \item \textit{Limite :} Performance souvent moindre sur des données très complexes et non structurées (images, texte).
    \end{itemize}
\end{frame}

\begin{frame}{Taxonomie des méthodes (2/2)}
    \textbf{2. Méthodes post-hoc (pour boîtes noires)}
    \vspace{0.3cm}
    Ces méthodes tentent d'expliquer un modèle complexe sans le modifier.
    \begin{itemize}
        \item \textbf{Indépendant du modèle (Model-Agnostic) :} Fonctionne sur n'importe quel type de modèle (LIME, SHAP).
        \item \textbf{Spécifique au modèle (Model-Specific) :} Utilise l'architecture interne spécifique (ex: Gradients ou cartes d'attention pour réseaux de neurones).
        \item \textbf{Explications Globales :} Comprendre le fonctionnement général du modèle.
        \item \textbf{Explications Locales :} Comprendre une décision spécifique pour un utilisateur donné.
    \end{itemize}
\end{frame}

% --- FIGURE VECTORIELLE LIME/SHAP ---
\begin{frame}{Illustration : approche locale (LIME et SHAP)}
    \begin{figure}
        \centering
        \begin{tikzpicture}[scale=0.7]
            % LIME
            \begin{scope}[shift={(-4,0)}]
                \draw[->] (-2,0) -- (2,0) node[right] {$x$};
                \draw[->] (0,-2) -- (0,2) node[above] {$y$};
                \draw[thick, blue] plot [smooth] coordinates {(-2,-1.5) (-0.5, 1.5) (1.5, 0.5)};
                \draw[red, dashed, thick] (-1.5, -0.5) -- (1, 2); 
                \fill[black] (-0.2, 0.8) circle (3pt) node[right] {Instance};
                \node at (0, -2.5) {\textbf{LIME} (Approx. Linéaire)};
            \end{scope}
            % SHAP
            \begin{scope}[shift={(4,0)}]
                \draw[->] (-1,0) -- (3,0) node[right] {Impact};
                \foreach \y/\w/\c/\l in {1.5/2/green/Age, 0.5/-1/red/Dette, -0.5/1/green/Revenu} {
                    \draw[fill=\c] (0,\y-0.2) rectangle (\w,\y+0.2);
                    \node[left] at (0,\y) {\l};
                }
                \draw (0,-1.5) -- (0,2);
                \node at (1, -2.5) {\textbf{SHAP} (Contributions)};
            \end{scope}
        \end{tikzpicture}
        \caption{Comparaison conceptuelle LIME vs SHAP (Généré par GEMINI PRO 3.0)}
    \end{figure}
\end{frame}

\begin{frame}{Focus méthode : LIME}
    \textbf{LIME (Local Interpretable Model-agnostic Explanations)}
    \begin{itemize}
        \item \textit{Principe :} LIME teste la sensibilité du modèle en perturbant légèrement l'entrée (ex: cacher des parties de l'image) et observe le changement de prédiction.
        \item \textit{Fonctionnement :} Il crée un modèle linéaire simple \textit{localement} autour de la prédiction pour l'approximer.
        \item \textit{Force :} Intuitif, visuel, applicable à tout modèle.
        \item \textit{Faiblesse :} Instabilité des explications (si on relance l'algorithme, l'explication peut changer légèrement).
    \end{itemize}
\end{frame}

\begin{frame}{LIME : fondement mathématique}
    \textbf{Explications locales interprétables indépendantes du modèle}
    
    \vspace{0.3cm}
    \textbf{Problème d'optimisation :}
    
    \[
    \text{explication}(\mathbf{x}) = \arg\min_{g \in G} \mathcal{L}(\hat{f}, g, \pi_{\mathbf{x}}) + \Omega(g)
    \]
    
    \vspace{0.2cm}
    où :
    \begin{itemize}
        \item \(\mathbf{x}\) : instance à expliquer
        \item \(\hat{f}\) : modèle boîte noire
        \item \(g\) : modèle simple interprétable (ex: régression linéaire)
        \item \(G\) : famille de modèles interprétables
        \item \(\mathcal{L}\) : fonction de perte pondérée par proximité
        \item \(\pi_{\mathbf{x}}\) : noyau de proximité (ex: \(\exp(-d^2/\sigma^2)\))
        \item \(\Omega(g)\) : mesure de complexité (régularisation)
    \end{itemize}
    
    \vspace{0.2cm}
    \textbf{Algorithme} : (1) Perturber \(\mathbf{x}\), (2) Prédire avec \(\hat{f}\), (3) Pondérer par \(\pi_{\mathbf{x}}\), (4) Ajuster \(g\) par régression pondérée.
\end{frame}

\begin{frame}{Focus méthode : SHAP}
    \textbf{SHAP (SHapley Additive exPlanations)}
    \begin{itemize}
        \item \textit{Principe :} Basé sur la théorie des jeux coopératifs (Valeurs de Shapley). Il calcule la contribution marginale de chaque caractéristique à la prédiction finale.
        \item \textit{Force :} Fondement mathématique solide, garantie de cohérence (propriété d'additivité).
        \item \textit{Faiblesse :} Très coûteux en temps de calcul pour des modèles complexes avec beaucoup de variables.
    \end{itemize}
\end{frame}

\begin{frame}{SHAP : fondement mathématique (valeurs de Shapley)}
    \textbf{Explications additives de Shapley}
    
    \vspace{0.3cm}
    \textbf{Définition (théorie des jeux coopératifs)} :
    
    \[
    \phi_j(val) = \sum_{S \subseteq \{1,\ldots,p\} \setminus \{j\}} \frac{|S|!(p-|S|-1)!}{p!} \left[ val(S \cup \{j\}) - val(S) \right]
    \]
    
    \vspace{0.2cm}
    où :
    \begin{itemize}
        \item \(\phi_j\) : contribution de la caractéristique \(j\) à la prédiction
        \item \(p\) : nombre total de caractéristiques
        \item \(S\) : sous-ensemble de caractéristiques (coalition)
        \item \(val(S)\) : prédiction du modèle en utilisant uniquement les caractéristiques de \(S\)
        \item \(\frac{|S|!(p-|S|-1)!}{p!}\) : poids combinatoire (contribution marginale moyenne)
    \end{itemize}
    
    \vspace{0.2cm}
    \textbf{Propriétés garanties} : Efficacité (somme = prédiction), Symétrie, Additivité, Joueur nul.
\end{frame}

\begin{frame}{Évaluation de l'explication}
    Comment savoir si une explication est "bonne" ? \parencite{Burkart2021A}
    \begin{itemize}
        \item \textbf{Fidélité (Fidelity) :} L'explication reflète-t-elle fidèlement le processus de décision du modèle ? (Mesure technique).
        \item \textbf{Robustesse (Robustness) :} Si l'entrée change très peu, l'explication reste-t-elle stable ?
        \item \textbf{Compréhensibilité (Understandability) :} L'explication est-elle compréhensible par l'humain cible ? (Mesure subjective/utilisateur).
        \item \textbf{Suffisance (Sufficiency) :} L'explication fournit-elle assez d'informations pour que l'humain puisse prédire le comportement du modèle ?
    \end{itemize}
\end{frame}

\begin{frame}{Évaluation de l'explication : critères formels}
    \textbf{Mesures d'évaluation des explications en IA Explicable} :
    
    \vspace{0.3cm}
    \begin{enumerate}
        \item \textbf{Fidélité} : 
        \[
        \text{Fidélité} = 1 - \frac{1}{N}\sum_{i=1}^{N}|\hat{f}(\mathbf{x}_i) - g(\mathbf{x}_i)|
        \]
        L'explication \(g\) reflète-t-elle \(\hat{f}\) localement ?
        
        \item \textbf{Robustesse} : Si \(\mathbf{x}' \approx \mathbf{x}\), alors \(\text{expl}(\mathbf{x}') \approx \text{expl}(\mathbf{x})\)
        
        \item \textbf{Compréhensibilité} : 
        \begin{itemize}
            \item Simplicité du modèle \(g\) (nombre de caractéristiques actives)
            \item Mesure subjective selon l'audience cible
        \end{itemize}
        
        \item \textbf{Complétude} : L'explication permet-elle de prédire le comportement de \(\hat{f}\) ?
    \end{enumerate}
    
    \vspace{0.2cm}
    \textbf{Compromis LIME vs SHAP} : LIME est rapide mais instable ; SHAP est lent mais théoriquement fondé.
\end{frame}

\begin{frame}{Les différentes audiences de l'explication}
    Une "bonne" explication dépend de son public \parencite{Doran2002What} :
    \begin{itemize}
        \item \textbf{Le Concepteur (Ingénieur) :} A besoin d'une explication technique pour déboguer (ex: valeurs SHAP, activation des neurones).
        \item \textbf{L'Utilisateur Final (Client) :} A besoin d'une explication simple, exploitable et non technique (ex: "Recommandé car vous avez aimé X et Y").
        \item \textbf{Le Régulateur (Auditeur) :} A besoin d'une preuve de non-discrimination et de robustesse (ex: analyse globale des biais).
    \end{itemize}
\end{frame}

\begin{frame}{Risques : le "pari de l'explicabilité"}
    L'explicabilité n'est pas une solution magique.
    \begin{itemize}
        \item \textbf{Fausse confiance (False Trust) :} Une explication plausible mais techniquement fausse peut amener un humain à faire trop confiance à un modèle défaillant.
        \item \textbf{Détournement du système (Gaming the System) :} Si un modèle est trop interprétable, des acteurs malveillants peuvent l'exploiter (ex: spam, fraude).
        \item \textbf{Rationalisation a posteriori :} L'explication peut n'être qu'une "bonne histoire" que le modèle raconte pour justifier une décision prise sur des corrélations fallacieuses.
    \end{itemize}
\end{frame}

\begin{frame}{Limitations techniques et conceptuelles}
    D'après \textcite{Doran2002What} et \textcite{Burkart2021A} :
    \begin{itemize}
        \item \textbf{Le compromis précision / interprétabilité :} Simplifier pour expliquer peut fausser la réalité du modèle. Une explication trop simple peut être trompeuse.
        \item \textbf{Biais de confirmation :} L'humain a tendance à accepter une explication qui confirme ses croyances, même si elle est techniquement fausse.
        \item \textbf{Fidélité (Fidelity gap) :} L'explication est-elle vraiment fidèle au fonctionnement du modèle, ou est-ce une rationalisation a posteriori ?
    \end{itemize}
\end{frame}

% ================= SECTION 4 : RECOMMANDATION =================
\section{Systèmes de Recommandation}

\begin{frame}{Section 4 : Systèmes de recommandation}
    \begin{center}
        \Large \textbf{L'architecture du choix}
    \end{center}
    \vspace{1cm}
    \begin{quote}
        « Expliquez ce qu'est un système de recommandation... »
    \end{quote}
\end{frame}

\begin{frame}{Définition et rôle}
    Selon \textcite{Ko2020A}, un système de recommandation (RS) est :
    \begin{itemize}
        \item Un système de filtrage de l'information visant à présenter les éléments les plus pertinents à un utilisateur.
        \item \textit{But :} Prédire la préférence ou la note qu'un utilisateur donnerait à un item qu'il ne connait pas encore.
        \item \textit{Enjeu :} Résoudre la surcharge informationnelle (Information Overload) sur le web.
    \end{itemize}
    Il ne s'agit pas seulement de suggérer, mais de \textit{personnaliser} l'expérience et d'orienter l'attention.
\end{frame}

% --- FIGURE VECTORIELLE RECO ---
\begin{frame}{Illustration : filtrage collaboratif vs contenu}
    \begin{figure}
        \centering
        \begin{tikzpicture}[scale=0.85]
            % CONTENT BASED
            \begin{scope}[shift={(-4,0)}]
                \node[font=\bfseries] at (0, 2.5) {Basé sur le contenu};
                \node[draw, circle, fill=orange!30] (u) at (0,1) {Utilisateur};
                \node[draw, rectangle, fill=green!20] (i1) at (-1,-1) {Item A};
                \node[draw, rectangle, fill=green!20] (i2) at (1,-1) {Item B};
                \draw[<->, dashed] (i1) -- (i2) node[midway, below, font=\tiny] {Attributs Similaires};
                \draw[->] (u) -- (i1) node[midway, left, font=\tiny] {Aime};
                \draw[->, red] (u) -- (i2) node[midway, right, font=\tiny] {Rec.};
            \end{scope}
            
            % COLLABORATIVE
            \begin{scope}[shift={(4,0)}]
                \node[font=\bfseries] at (0, 2.5) {Filtrage collaboratif};
                \draw[step=0.5, gray!50] (-1.5,-1.5) grid (1.5,1.5);
                \node at (0, 1.8) {\scriptsize Matrice User-Item};
                \fill[blue] (-1, 1) circle (3pt);
                \fill[blue] (-0.5, 1) circle (3pt);
                \fill[red] (0, 0.5) circle (3pt) node[right, red, font=\tiny] {Prédiction};
            \end{scope}
        \end{tikzpicture}
        \caption{Distinctions conceptuelles FC vs FBC (Généré par GEMINI PRO 3.0)}
    \end{figure}
\end{frame}

\begin{frame}{Aperçu des techniques}
    Les principales familles identifiées par \textcite{Ko2020A} :
    \begin{enumerate}
        \item \textbf{Filtrage basé sur le contenu (Content-Based) :} Analyse les propriétés des objets. "Vous avez aimé Star Wars (Science-Fiction), voici Dune".
        \item \textbf{Filtrage collaboratif (Collaborative Filtering) :} Analyse les interactions sociales. "Les gens comme vous ont aimé ceci".
        \item \textbf{Systèmes hybrides :} Combinaison des deux approches pour pallier leurs défauts respectifs.
        \item \textbf{Basé sur l'apprentissage profond (Deep Learning) :} Utilisation de réseaux de neurones (RNN, CNN) pour modéliser des interactions complexes et séquentielles.
    \end{enumerate}
\end{frame}

% --- SLIDE MATHÉMATIQUE 1 : CB ---
\begin{frame}{Fondement mathématique (1) : Basé sur le contenu (FBC)}
    Dans l'approche FBC, la similarité est souvent calculée via la \textbf{similarité cosinus}.
    
    Soit \(u\) un vecteur de profil utilisateur et \(i\) un vecteur de profil d'item (basé sur TF-IDF par exemple). La pertinence est estimée par :
    
    \begin{equation}
        \text{sim}(u, i) = \cos(\theta) = \frac{u \cdot i}{||u|| \times ||i||} = \frac{\sum_{k=1}^{n} u_k i_k}{\sqrt{\sum_{k=1}^{n} u_k^2} \sqrt{\sum_{k=1}^{n} i_k^2}}
    \end{equation}
    
    \begin{itemize}
        \item \textbf{Principe :} On projette l'utilisateur et l'item dans le même espace vectoriel.
        \item Plus l'angle \(\theta\) est petit, plus le cosinus est proche de 1, signifiant une forte similarité.
    \end{itemize}
\end{frame}

\begin{frame}{Technique 1 : le filtrage collaboratif (FC)}
    Détaillé par \textcite{Koren2022Advances}, c'est la méthode la plus célèbre.
    \begin{itemize}
        \item \textbf{Principe :} Utiliser les interactions passées (matrice Utilisateurs-Items) pour prédire les interactions futures, sans connaitre la nature des objets.
        \item \textbf{Approche centrée utilisateur (User-based) :} Trouver des voisins similaires à l'utilisateur cible et recommander ce qu'ils ont aimé.
        \item \textbf{Approche centrée item (Item-based) :} Trouver des items qui sont notés de manière similaire par l'ensemble des utilisateurs.
        \item \textbf{Factorisation Matricielle (SVD) :} Décomposer la matrice géante des notations en facteurs latents (goûts cachés) pour remplir les vides.
    \end{itemize}
\end{frame}

% --- SLIDE MATHÉMATIQUE 2 : CF ---
\begin{frame}{Fondement mathématique (2) : factorisation matricielle (DVS)}
    Selon \textcite{Koren2022Advances}, le but est d'approximer la matrice de notation \(R\) par deux matrices de facteurs latents \(P\) (utilisateurs) et \(Q\) (items) de dimension \(f\).
    
    L'estimation de la note \(\hat{r}_{ui}\) est donnée par le produit scalaire :
    \[ \hat{r}_{ui} = q_i^T p_u \]
    
    Pour apprendre ces facteurs, on minimise l'erreur quadratique régularisée sur l'ensemble des notes connues \(\kappa\) :
    
    \begin{equation}
        \min_{q^*, p^*} \sum_{(u,i) \in \kappa} (r_{ui} - q_i^T p_u)^2 + \lambda(||q_i||^2 + ||p_u||^2)
    \end{equation}
    
    \footnotesize
    \textit{Le terme \(\lambda\) (régularisation) empêche le surapprentissage (overfitting) en pénalisant les vecteurs de trop grande amplitude.}
\end{frame}

\begin{frame}{Technique 2 : filtrage basé sur le contenu (FBC)}
    \begin{itemize}
        \item \textbf{Principe :} Créer un profil pour chaque item (caractéristiques : genre, auteur, mots-clés) et un profil utilisateur basé sur ses préférences déclarées ou son historique.
        \item \textbf{Méthode :} Calcul de similarité (Distance Cosinus, Jaccard) entre le vecteur Utilisateur et le vecteur Item.
        \item \textbf{Avantage :} Indépendance. Pas besoin des données des autres utilisateurs.
        \item \textbf{Limite :} Sur-spécialisation. Le système ne recommande que ce que l'utilisateur connait déjà, limitant la découverte (sérendipité).
    \end{itemize}
\end{frame}

\begin{frame}{Le problème du "démarrage à froid" (Cold Start)}
    Un défi majeur pour les systèmes collaboratifs \parencite{Ko2020A} :
    \begin{itemize}
        \item \textbf{Nouvel utilisateur :} Comment recommander à quelqu'un dont on ne sait rien ? Le système n'a pas de "voisins" à qui le comparer.
        \item \textbf{Nouvel item :} Comment recommander un nouveau film si personne ne l'a encore noté ? Il n'apparaîtra dans les recommandations de personne.
        \item \textbf{Solution hybride :} C'est là que le FBC est crucial. On peut recommander un nouvel item en se basant sur ses attributs (genre, acteurs) avant même qu'il ait reçu des notes.
    \end{itemize}
\end{frame}

\begin{frame}{Évaluation des recommandations (métriques)}
    Comment mesurer la qualité d'un SR ?
    \begin{itemize}
        \item \textbf{Précision (Precision@k) :} Parmi les k items recommandés, combien sont pertinents ?
        \item \textbf{Rappel (Recall@k) :} Parmi tous les items pertinents, combien ont été recommandés dans le top k ?
        \item \textbf{Métriques de classement (NDCG) :} La pertinence seule ne suffit pas. Les meilleurs items doivent-ils être en position 1 ou 2 ?
        \item \textbf{Métriques "au-delà de la précision" :}
            \begin{itemize}
                \item \textbf{Sérendipité :} Capacité à surprendre l'utilisateur avec des items pertinents mais inattendus.
                \item \textbf{Diversité :} Les recommandations couvrent-elles différents aspects des goûts de l'utilisateur ?
            \end{itemize}
    \end{itemize}
\end{frame}

\begin{frame}{Mesures d'évaluation : formules mathématiques}
    \textbf{Mesures de classement pour systèmes de recommandation}
    
    \vspace{0.3cm}
    \textbf{1. Précision@k} :
    \[
    \text{Précision@}k = \frac{|\{\text{items pertinents}\} \cap \{\text{top-}k\text{ recommandés}\}|}{k}
    \]
    
    \textbf{2. Rappel@k} :
    \[
    \text{Rappel@}k = \frac{|\{\text{items pertinents}\} \cap \{\text{top-}k\text{ recommandés}\}|}{|\{\text{tous items pertinents}\}|}
    \]
    
    \textbf{3. GCDAN@k (Gain Cumulé Actualisé Normalisé)} :
    \[
    \text{GCDAN@}k = \frac{\text{GCA@}k}{\text{GCAI@}k}, \quad \text{GCA@}k = \sum_{i=1}^{k} \frac{2^{rel_i} - 1}{\log_2(i+1)}
    \]
    
    Prend en compte l'ordre et la pertinence graduée.
    
    \vspace{0.2cm}
    \textbf{Au-delà de la précision} : Diversité, Sérendipité, Nouveauté, Couverture.
\end{frame}

\begin{frame}{Enjeux : la "bulle de filtre" et sérendipité}
    Le SR façonne notre vision du monde.
    \begin{itemize}
        \item \textbf{Bulle de filtre :} L'optimisation de la pertinence (clics) peut enfermer l'utilisateur dans une boucle de rétroaction, ne lui montrant que ce qu'il pense déjà.
        \item \textbf{Perte de sérendipité :} Un système trop efficace en FC ou FBC ne laisse plus de place à la découverte fortuite, essentielle à l'apprentissage et à la culture.
        \item \textbf{Enjeu éthique :} Le SR a une responsabilité dans la polarisation du débat public (Q2) et devient une forme de bureaucratie de l'information (Q1).
    \end{itemize}
\end{frame}

\begin{frame}{Enjeux éthiques : bulle de filtre et polarisation}
    \textbf{Bulle de filtre et chambres d'écho}
    
    \vspace{0.3cm}
    \textbf{Définitions formelles} :
    \begin{itemize}
        \item \textbf{Bulle de filtre (Pariser, 2011)} : Isolation intellectuelle causée par personnalisation algorithmique sélective
        \item \textbf{Chambre d'écho} : Renforcement des croyances par exposition homogène (boucles de rétroaction)
    \end{itemize}
    
    \vspace{0.3cm}
    \textbf{Mécanismes algorithmiques} :
    \begin{enumerate}
        \item FC/FBC favorisent la similarité $\rightarrow$ Recommandations redondantes
        \item Boucle de rétroaction : \(\text{Préférence}_t \rightarrow \text{Reco}_t \rightarrow \text{Interaction}_t \rightarrow \text{Préférence}_{t+1}\)
        \item Optimisation pour engagement (temps passé) \(\neq\) diversité informationnelle
    \end{enumerate}
    
    \vspace{0.3cm}
    \textbf{Atténuation} : Injection de diversité, "données antidotes", mesures au-delà de la précision (sérendipité, diversité).
\end{frame}

\begin{frame}{Comparaison détaillée}
    \begin{table}[]
        \centering
        \small
        \begin{tabular}{@{}p{5cm}p{5cm}@{}}
            \toprule
            \textbf{Filtrage collaboratif (FC)} & \textbf{Basé sur le contenu (FBC)} \\ \midrule
            Nécessite une communauté (données sociales) & Autonome (centré sur l'utilisateur seul) \\
            Favorise la \textit{Sérendipité} (peut surprendre) & Tendance à la sur-spécialisation (bulle) \\
            Problème de "Démarrage à froid" (nouveaux items/utilisateurs) & Gère bien les nouveaux items \\
            Opaque (facteurs latents mathématiques) & Plus explicable (basé sur des attributs clairs) \\ \bottomrule
        \end{tabular}
    \end{table}
    \textcite{Ko2020A} notent que les systèmes modernes (Netflix, YouTube) sont presque toujours hybrides pour maximiser les forces des deux approches.
\end{frame}

\begin{frame}{L'évolution vers l'Apprentissage Profond}
    Les limites du FC (linéarité) sont dépassées par l'Apprentissage Profond (Deep Learning).
    \begin{itemize}
        \item \textbf{Filtrage Collaboratif Neuronal (FCN) :} Remplace le simple produit scalaire par un réseau de neurones pour capturer des interactions complexes.
        \item \textbf{Réseaux de Neurones Récurrents (RNN) :} Utilisés pour la recommandation de session (séquentielle), prédisant le prochain clic en fonction de l'historique immédiat.
        \item \textbf{Enjeu :} Augmentation massive de l'opacité. On gagne en précision, mais on perd en capacité d'explication \(\rightarrow\) Retour à la question 3 (Interprétabilité).
    \end{itemize}
\end{frame}

\begin{frame}{Filtrage collaboratif neuronal (FCN)}
    \textbf{Architecture d'apprentissage profond pour la recommandation}
    
    \vspace{0.3cm}
    \textbf{Limite de la FM traditionnelle} : Le produit scalaire \(\mathbf{u}_u^T \cdot \mathbf{v}_i\) est linéaire et ne capture pas les interactions complexes.
    
    \vspace{0.3cm}
    \textbf{FCN : Remplacer le produit scalaire par un réseau de neurones} :
    
    \[
    \hat{y}_{ui} = f(\mathbf{u}_u, \mathbf{v}_i | \Theta)
    \]
    
    où \(f\) est un Perceptron Multi-Couches (PMC).
    
    \vspace{0.3cm}
    \textbf{Architecture FMN (Factorisation Matricielle Neuronale)} :
    \begin{enumerate}
        \item \textbf{FMG (FM Généralisée)} : Produit élément-par-élément des plongements
        \item \textbf{PMC} : Couches denses empilées (4-5 couches, non-linéarité ReLU)
        \item \textbf{Fusion} : Concaténation FMG + PMC $\rightarrow$ couche de sortie
    \end{enumerate}
    
    \[
    \hat{y}_{ui} = \sigma\left( \mathbf{h}^T \left[ \mathbf{u}_u \odot \mathbf{v}_i \, || \, \text{PMC}(\mathbf{u}_u, \mathbf{v}_i) \right] \right)
    \]
    
    \textbf{Gain} : Capture d'interactions non-linéaires.
\end{frame}

% ================= CONCLUSION =================
\section{Conclusion}

% --- SYNTHÈSE TRANSVERSALE (La diapo qui manquait) ---
\begin{frame}{Synthèse transversale : architecture du Léviathan}
    \textbf{Interconnexions entre les 4 questions} :
    
    \vspace{0.5cm}
    \begin{center}
    \begin{tikzpicture}[node distance=2.5cm, auto, box/.style={rectangle, draw, rounded corners, fill=gray!10, align=center, minimum height=1cm}]
        % Noeuds
        \node[box, fill=blue!20] (Q1) {\textbf{Q1 : Bureaucratie}\\ (La Logique)};
        \node[box, fill=green!20, right=of Q1] (Q4) {\textbf{Q4 : Recommandation}\\ (Le Moteur)};
        \node[box, fill=red!20, below=of Q1] (Q2) {\textbf{Q2 : Déontologie}\\ (La Loi)};
        \node[box, fill=yellow!20, right=of Q2] (Q3) {\textbf{Q3 : IAX}\\ (Le Contre-Pouvoir)};
        
        % Flèches
        \draw[->, thick] (Q1) -- node[above] {\footnotesize Implémente} (Q4);
        \draw[->, thick] (Q4) -- node[left] {\footnotesize Crée opacité} (Q3);
        \draw[->, thick] (Q2) -- node[right] {\footnotesize Exige} (Q3);
        \draw[<->, dashed, thick] (Q1) -- node[left] {\footnotesize Encadre} (Q2);
        
    \end{tikzpicture}
    \end{center}
    
    \vspace{0.3cm}
    \footnotesize
    \begin{enumerate}
        \item La \textbf{bureaucratie algorithmique} s'incarne techniquement dans les \textbf{systèmes de recommandation}.
        \item L'opacité de ces systèmes (boîte noire) rend nécessaire l'\textbf{IA explicable}.
        \item L'ensemble doit être régulé par une \textbf{déontologie} forte pour éviter la dérive (capitalisme de surveillance).
    \end{enumerate}
\end{frame}

\begin{frame}{Synthèse finale}
    \begin{itemize}
        \item L'algorithme de personnalisation est le \textbf{bureaucrate parfait} de Weber, efficace mais aveugle, appliquant des règles impersonnelles à grande échelle.
        \item Face à ce pouvoir, un \textbf{code de déontologie} inspiré de la Biosphère est nécessaire pour responsabiliser les créateurs et sortir du vide éthique.
        \item Cette responsabilité passe techniquement par l'\textbf{interprétabilité (XAI)}, seule garante contre l'arbitraire et condition de la confiance.
        \item Les \textbf{systèmes de recommandation} sont le terrain d'application concret de ces tensions, tiraillés entre la performance pure (Apprentissage profond) et la nécessité de transparence et d'équité.
    \end{itemize}
\end{frame}

\begin{frame}{Conclusion}
    \begin{center}
        \Huge Le Léviathan peut être dompté, \\ mais seulement s'il est compris.
    \end{center}
\end{frame}

% --- GLOSSAIRE ---
\begin{frame}{Glossaire : terminologie francisée}
    \footnotesize
    \begin{table}[]
        \centering
        \begin{tabular}{@{}p{5cm}p{6cm}@{}}
            \toprule
            \textbf{Terme anglais} & \textbf{Terme français} \\ \midrule
            Digital Weberianism & Wébérisme numérique \\
            Algorithmic Accountability & Redevabilité algorithmique \\
            XAI (Explainable AI) & IA Explicable / IAX \\
            LIME & Explications Locales Interprétables Indépendantes du Modèle \\
            SHAP & Explications Additives de Shapley \\
            SVD & DVS (Décomposition en Valeurs Singulières) \\
            TF-IDF & FT-FDI (Fréquence du Terme - Fréquence du Document Inverse) \\
            Matrix Factorization (MF) & Factorisation Matricielle (FM) \\
            Neural Collaborative Filtering (NCF) & Filtrage Collaboratif Neuronal (FCN) \\
            NeuMF & FMN (Factorisation Matricielle Neuronale) \\
            MLP (Multi-Layer Perceptron) & PMC (Perceptron Multi-Couches) \\
            GMF (Generalized MF) & FMG (FM Généralisée) \\
            NDCG & GCDAN (Gain Cumulé Actualisé Normalisé) \\
            Precision@k / Recall@k & Précision@k / Rappel@k \\
            Collaborative Filtering (CF) & Filtrage Collaboratif (FC) \\
            Content-Based (CB) & Basé sur le Contenu (FBC) \\
            Recommender System (RS) & Système de Recommandation (SR) \\
            Filter Bubble / Echo Chamber & Bulle de Filtre / Chambre d'Écho \\
            Information Overload & Surcharge informationnelle \\
            Gaming the System & Détournement du système \\
            Model-Agnostic & Indépendant du modèle \\
            Model-Specific & Spécifique au modèle \\
            Features & Caractéristiques / Attributs \\
            Cold Start & Démarrage à froid \\
            AI Act & Loi sur l'IA \\ \bottomrule
        \end{tabular}
    \end{table}
\end{frame}

% --- BIBLIOGRAPHIE ---
\begin{frame}[allowframebreaks]{Bibliographie}
    \printbibliography
\end{frame}

\end{document}